\documentclass[letterpaper,10pt]{article}
\usepackage[utf8]{inputenc}
\usepackage[T1]{fontenc}
\usepackage[spanish]{babel}
\usepackage{hyperref}
\usepackage{geometry}
\geometry{margin=0.85in}
\usepackage{enumitem}
\setlist[itemize]{noitemsep, topsep=2pt, leftmargin=*}

\begin{document}

\begin{center}
    {\Large \textbf{Fernando Rodríguez Prianti}}\\[0.1cm]
    {\large Ingeniero de Software · Full Stack}\\[0.15cm]
    Xalapa, Veracruz, México \\
    \href{mailto:fernando.rdz.prianti@gmail.com}{fernando.rdz.prianti@gmail.com} \quad | \quad +52 228 537 2622 \\
    \href{https://www.linkedin.com/in/fernando-rodr{\'i}guez-732670140/}{LinkedIn} \quad | \quad
    \href{https://github.com/FernandoRP-nan}{GitHub} \quad | \quad
    \href{https://fernandorp-nan.github.io/fernando-cv/}{CV Web}
\end{center}

\section*{Perfil}
Ingeniero de Software con 4+ años en web y móvil. PWAs y portales en producción (Angular, Laravel), apps Android con Jetpack Compose y ecosistemas multi-repo. Experiencia en Stripe, JWT/RBAC, push FCM, WebSockets y Docker. Side projects: Python, Godot, plugins Obsidian y RAG local (LangGraph + Ollama).

\section*{Experiencia}
\textbf{Tierra de Perros} \hfill \textit{DIC 2024 -- DIC 2025 · 13 meses}\\
Xalapa, Veracruz
\begin{itemize}
    \item PWA Angular + Laravel: citas, tienda, Stripe, chat WebSocket y roles.
    \item Push FCM, JWT. Plataforma en operación; contrato finalizado dic 2025.
\end{itemize}

\textbf{Portal Bachilleres Veracruz} \hfill \textit{SEP 2024 -- DIC 2024 · 4 meses}\\
Xalapa, Veracruz
\begin{itemize}
    \item Portal gubernamental PHP + JS: noticias, horarios, avisos y JWT por rol.
\end{itemize}

\textbf{Satori Tech -- Desarrollador Android Jr.} \hfill \textit{SEP 2022 -- JUN 2024 · 22 meses}\\
Mecatepec, CDMX
\begin{itemize}
    \item Apps Kotlin con MVVM, Retrofit, Room y metodología ágil.
\end{itemize}

\textbf{Grupo Modelo -- Consultor SAP PI PO Jr.} \hfill \textit{JUN -- AGO 2022 · 3 meses}\\
Xalapa, Veracruz
\begin{itemize}
    \item Soporte SAP ERP.
\end{itemize}

\section*{Proyectos destacados}
\textbf{Pokédex} (2024 · 4 meses): Kotlin, Compose, Room, Retrofit, MVVM.\\
\textbf{Uno Bank} (2026): Compose, Clean Architecture, Paging, Hilt, biometría, backend simulado.\\
\textbf{Cero — Películas \& Compose} (2026): Compose, Hilt, Room, Retrofit, TMDB, módulo NDK.\\
\textbf{Punto de Venta} (2024): Node.js, React, MySQL, JWT, roles gerente/cajero.\\
\textbf{Gallery AYN} (2024): Python, Svelte, escritorio con galería y video.\\
\textbf{All You Need} (2022--2025): Ecosistema Android Compose + Spring Boot + CouchDB.\\
\textbf{Buscador DFS} (2024): Python, 8-puzzle con DFS y validación de estados.\\
\textbf{Vault Plugins} (2024--2025): TypeScript, Obsidian API (tareas, colección, agenda).\\
\textbf{IA Obsidian Manager} (2025): LangGraph, Ollama, ChromaDB, Docker.

\section*{Stack técnico}
\textbf{Lenguajes:} Java, Kotlin, Python, C\#, JavaScript, TypeScript, PHP, GDScript.\\
\textbf{Frameworks:} Angular, Laravel, React, Node.js, Express, Jetpack Compose, Spring Boot, Svelte, Godot.\\
\textbf{Bases de datos:} MySQL, PostgreSQL, SQLite, MongoDB, Room, CouchDB.\\
\textbf{Integraciones:} Stripe, Firebase FCM, WebSockets, REST API, JWT.\\
\textbf{Herramientas:} Git, GitHub Actions, Docker, Android Studio, VS Code, IntelliJ, Vite, Gradle.\\
\textbf{Prácticas:} MVVM, Clean Architecture, Repository, inyección de dependencias, patrones de diseño.\\
\textbf{IA:} LangGraph, Ollama, ChromaDB, RAG.

\section*{Educación}
\textbf{Ingeniería en Sistemas Computacionales -- Especialización en Software}\\
Instituto Tecnológico Superior de Xalapa \hfill 2019 -- 2025

\section*{Idiomas}
Español (nativo) · Inglés (A2)

\end{document}
